\chapter{Validazione empirica}

\section{Dataset}

Per la validazione sperimentale è stato utilizzato il
\textit{PE Malware Machine Learning Dataset} di Practical Security
Analytics \cite{lester2021}. Il dataset comprende, oltre ai file PE,
un file CSV contenente diverse informazioni associate a ciascun
file, tra cui la lunghezza in byte e l'etichetta che indica
l'appartenenza alla \textit{Whitelist} o alla \textit{Blacklist},
permettendo di distinguere rispettivamente i file benigni da quelli
malevoli. Queste caratteristiche rendono il dataset particolarmente
adatto al lavoro svolto, poiché mette a disposizione sia i file PE
originali, necessari per la loro conversione in immagini, sia le
relative etichette utilizzate per la classificazione.

Il dataset è composto complessivamente da 201.549 file PE, di cui
86.812 appartenenti alla \textit{Whitelist} e quindi considerati
benigni, e 114.737 appartenenti alla \textit{Blacklist} e quindi
considerati malevoli.

Ai fini della validazione sperimentale, il dataset è stato suddiviso
in un training set, contenente l'80\% dei campioni, e un test set,
contenente il restante 20\%. La suddivisione è stata effettuata in
modo stratificato rispetto alla classe, così da mantenere la
proporzione tra file benigni e malevoli presente nel dataset
originale.

\begin{table}[!ht]
    \centering
    \caption{Distribuzione dei campioni nei set di training e test}
    \label{tab:dataset_split}
    \renewcommand{\arraystretch}{1.3}
    \begin{tabular}{lrrr}
        \toprule
        \textbf{Insieme} & \textbf{Benigni} & \textbf{Malevoli}
        & \textbf{Totale} \\
        \midrule
        Training & 69.450 & 91.789 & 161.239 \\
        Test     & 17.362 & 22.948 & 40.310 \\
        \midrule
        Totale   & 86.812 & 114.737 & 201.549 \\
        \bottomrule
    \end{tabular}
\end{table}

Il training set è stato utilizzato per l'addestramento del modello,
mentre il test set è stato utilizzato per valutarne le prestazioni,
dato che contiene dati non utilizzati durante il training e quindi
permette di valutare il modello su campioni mai utilizzati durante
l'addestramento.

\section{Set-up e metriche}



\section{Risultati}